% Clase 19 — planteo del desarrollo multipolar magnético. El origen se toma
% cerca del circuito, de modo que |r_1| es del orden de su tamaño a, y se mira
% el campo en un punto lejano, |r_2| >> a.
\begin{tikzpicture}[scale=1]

  % circuito
  \draw[figblue, line width=1.3pt]
    plot[smooth cycle, tension=0.8] coordinates {(-1.0,-0.1) (-0.5,0.85) (0.6,1.0) (1.2,0.2) (0.7,-0.8) (-0.4,-0.85)};
  \draw[figvecres] (-0.05,-0.97) -- (0.5,-0.9);
  \node[figlbl, text=figred] at (0.15,-1.3) {$I$};
  \node[figlbl] at (-1.25,0.75) {$C$};

  % origen y r_1
  \fill[figgray] (0,0) circle (1.4pt);
  \node[figlblaux] at (-0.22,-0.22) {$O$};
  \draw[figvec] (0,0) -- (0.5,0.97);
  \node[figlbl] at (-0.08,0.6) {$\vec r_1$};

  % tamaño a
  \draw[figgray, line width=0.6pt, {Latex[length=1.6mm]}-{Latex[length=1.6mm]}] (-1.05,-1.75) -- (1.25,-1.75);
  \node[figlblaux] at (0.1,-2.05) {$\sim a$};

  % punto lejano
  \coordinate (P) at (8.2,2.1);
  \draw[figvec] (0,0) -- (P);
  \node[figlbl] at (4.4,0.75) {$\vec r_2$};
  \draw[figaux, -{Latex[length=2mm,width=1.4mm]}] (0.5,0.97) -- (P);
  \node[figlbl] at (3.6,1.9) {$\vec r_2-\vec r_1$};
  \fill[figblue] (P) circle (1.6pt);
  \node[figlbl, anchor=west] at (7.0,2.6) {$|\vec r_2| \gg a$};

\end{tikzpicture}
